\section{An alternating extraction layer excludes a large triangle partner}
\label{sec:standard-triangle-alternating-layer}

\paragraph{Status and shared provenance.}
D-TD46, reserved question Q2fb, is the triangle refinement following
D-TD44. Sixfolds first considered a second stable standard layer;
root proposed the stronger use of any alternating extraction layer,
the prime-order deficit estimate and the constant thirty. Both
independently reconstructed the actual local-normalizer interface.
The complete proof is written below by \texttt{/root/sixfolds}.
Complete author whole-proof, source-interface, checker and first-output
readings pass. Independent and root whole-proof reviews also pass.
Imprimitive's separately assigned D-TD47 generalizes the deficit step;
no independent movement is claimed twice for overlapping consequences.

Retain the actual sphere-source setting of HD-SS08--11:
\[
 \begin{gathered}
 V=(\operatorname{Std}_n\circ x)\otimes U,\qquad
 x:J\twoheadrightarrow T\in\{A_n,S_n\},\quad n\ge5,\\
 b=\dim U\ge2,\qquad \mu_V=6,\qquad
 K_0=\ker U,\qquad Q=U(J),\qquad \mu=\mu_U>0.
 \end{gathered}
 \tag{SAT.1}
\]
Here V is faithful and absolutely irreducible, both tensor factors
are actual finite linear representations, and the connected
J-Galois source C has base \(\mathbf P^1\). The coefficient V has
a model over its exact imaginary quadratic character field K.
Assume \(K_0\) is transitive on the n natural letters. Put
\[
 H=x^{-1}((S_n)_1),\quad N=H\cap K_0,\quad
 X=C/H,\quad Y=C/N,\quad C_U=C/K_0,\quad h=g(X).
 \tag{SAT.2}
\]
The actual Q-cover \(C_U\to\mathbf P^1\) is assumed to have
exact signature \((2,3,7)\). Its three generators will be denoted
\(q_2,q_3,q_7\), with the indicated exact orders and product one.

\subsection*{The actual alternating-layer hypothesis}

Suppose the entire U contains one actual stable extracted matrix
algebra on a space B. Its projective image P is an actual quotient
of Q, with the full normal subgroup and irreducible tensor action
\[
 A_m^t\triangleleft P\le\operatorname{Aut}(A_m^t),\qquad
 m\ge8,\quad t\ge1,\quad
 B=W_1\otimes\cdots\otimes W_t,\quad
 U=B\otimes R.
 \tag{SAT.3}
\]
The individual component algebras are transitively permuted by P;
their finite perfect central lifts act irreducibly and nontrivially
on the equal-dimensional spaces \(W_j\). The entire external
partner R remains present. These are precisely the nonabelian
coarse extraction properties of D-TD12 and D-TD17--18 when the
simple group is \(A_m\). No individual component is declared
Q-stable. No represented perfect lift is asserted to be a normal
subgroup of the original Q or J.

For \(m\ge8\), the natural description
\(\operatorname{Aut}(A_m)=S_m\) gives an actual action
\[
 Q\longrightarrow P\longrightarrow S_m\wr S_t
       \longrightarrow S_{mt}.
 \tag{SAT.4}
\]
It acts on t disjoint natural m-letter sets. The full normal
\(A_m^t\) is transitive on each set, and P is transitive on the
sets, so this degree \(D=mt\) action is transitive. This is an
actual permutation action of Q, rather than an alternating
composition factor or a natural action chosen independently at
each branch point.

\begin{theorem}[D-TD46]
Under (SAT.1)--(SAT.4), if \(D=mt\ge27\), then
\[
                  b\ge3,\qquad n\le4596.
 \tag{SAT.5}
\]
This conclusion uses neither period rank nine nor source minimality.
It permits every irreducible component module and any number of
permuted components in (SAT.3).

In the nonfaithful-partner, primitive, minimal exact-K, rank-nine,
naturally polarized setting of D-TD44, with the same fixed complete
extraction decomposition and
selected standard layer, the actual partner-cover signature
\((2,3,7)\) therefore forces
\[
                              n\le8628.
 \tag{SAT.6}
\]
Here minimality is among actual sources of the same exact field,
signature \((3,3)\), actual period rank nine, and natural signed
polarization determinant class, exactly as in HD-SS06 and D-TD44.
No statement about a source assigned separately to U is made.
\end{theorem}

\begin{proof}
\paragraph{Prime-order deficits in the whole natural action.}
Let \(z_d\in S_D\) be the action of \(q_d\) in (SAT.4),
for \(d=2,3,7\). Its order divides d. Since d is prime, all
cycles have length one or d, including the possible identity case.
If \(s_d\) is its number of moved letters, then
\[
 \operatorname{ind}(z_d)=(1-1/d)s_d,\qquad
 \delta_d=(1-1/d)D-\operatorname{ind}(z_d)
           =(1-1/d)(D-s_d)\ge0.
 \tag{SAT.7}
\]
The transitive Q-set (SAT.4) defines an actual connected degree-D
quotient curve Z of the given Q-cover. Riemann--Hurwitz on this
curve, with identity branch images allowed, gives
\[
 \begin{split}
 \sum_{d=2,3,7}\operatorname{ind}(z_d)&=2D+2g(Z)-2,\\
 \sum_{d=2,3,7}\delta_d&=D/42+2-2g(Z)\le D/42+2.
 \end{split}
 \tag{SAT.8}
\]
Every individual nonnegative deficit is at most this sum.
Consequently
\[
 s_d\ge D-\frac{d}{d-1}(D/42+2)
       \ge\frac{20D}{21}-4.
 \tag{SAT.9}
\]
Every nontrivial power \(q_d^a\), \(1\le a<d\), has the same
cycle lengths and fixed points as \(z_d\) in the natural action,
so (SAT.9) holds for
all those powers. This statement concerns powers of the actual
Q-inertia elements; it does not replace the periods of the
original coefficient operators. For \(D\ge27\),
\[
             \frac{20}{21}-\frac4D\ge\frac45,
       \qquad\frac1{24}\left(\frac{20}{21}-\frac4D\right)
                      \ge\frac1{30}.
 \tag{SAT.10}
\]
In particular none of the indicated natural images is the identity.

\paragraph{Actual finite local normalizers and scalar phases.}
Fix one of the powers \(q=q_d^a\) just considered. If it permutes
any component algebras in (SAT.3), D-TD14 gives
\(\nu_U(U(q))/b\ge1/4\), retaining every other tensor factor.
Here \(\nu\) denotes the codimension of the largest eigenspace,
independently of its eigenvalue. Thus this case already exceeds
the required bound \(1/30\).

Otherwise q fixes every component algebra. Its natural action on
the D letters is the disjoint union of \(y_j\in S_m\).
Their supports sum to at least the right side of (SAT.9), so some
component satisfies
\[
 \frac{k_j}{m}:=\frac{\operatorname{mov}(y_j)}m
                 \ge\frac{20}{21}-\frac4{mt}>0.
 \tag{SAT.11}
\]
We verify the exact central-cover hypotheses for that component.
Let E be its actual finite perfect central lift of \(A_m\),
acting irreducibly on \(W_j\), and let A be a local matrix
representing q. Its projective class belongs to the finite actual
component-stabilizer image and normalizes the projective E.
Thus A normalizes \(\mathbf C^*E\), and perfectness gives
\([\mathbf C^*E,\mathbf C^*E]=E\). Conjugating this equality
shows that A honestly normalizes E. This proves normalization
without assuming a lift of E inside the original source group.

The projective order of A is finite, so a scalar rescaling gives
\(A^L=1\) for some positive L. This does not change any largest
eigenspace dimension. The group \(G=\langle E,A\rangle\) is
finite: E is finite and normal and G is a union of at most L
cosets of E. Conjugation gives
\(G\to\operatorname{Aut}(A_m)=S_m\). Its kernel is scalar.
Indeed an element in the kernel conjugates each \(e\in E\)
to \(c(e)e\); the scalar function c is a homomorphism and is
trivial by perfectness. The element centralizes E and is scalar
by irreducibility and Schur's lemma. This kernel is central.

If \(y_j\) is even, choose \(e\in E\) with the same inner
automorphism. Then \(Ae^{-1}\) belongs to that scalar kernel.
The even assertion of D-TD31 applies to e and hence to A.
If \(y_j\) is odd, the image of G is all of \(S_m\), its
kernel is central, G is generated by E and A, and restriction
to E is irreducible and nontrivial. These are exactly the
explicit odd-extension hypotheses of D-TD31. In either case,
\[
             \frac{\nu_{W_j}(A)}{\dim W_j}
                         \ge\frac{k_j}{24m}.
 \tag{SAT.12}
\]
The extraction field or an abstract alternating quotient is not
substituted for this actual finite local linear group.

Since q fixes all component algebras, its whole operator on U is
a tensor product of local matrices and the entire partner matrix,
with arbitrary reciprocal scalar choices. The whole operator is
finite order; each factor has a scalar power and is diagonalizable
over \(\mathbf C\). Spectral convolution cannot increase the
largest normalized eigenvalue multiplicity of the selected factor.
Combining (SAT.10)--(SAT.12), including the permuted case, proves
\[
 \frac{\operatorname{rank}(1-U(q_d^a))}{b}
       \ge\frac{\nu_U(U(q_d^a))}{b}\ge\frac1{30},
 \qquad d=2,3,7,\quad 1\le a<d.
 \tag{SAT.13}
\]
This is a bound on this complete set of actual inertia powers,
not a new global support floor for every operator in Q.

\paragraph{Connected base change and all descended inertia.}
Transitivity of \(K_0\) gives \(J=K_0H\) and
\(H/N\simeq Q\). The fixed-field compositum identifies Y
with the connected normalization of
\(X\times_{\mathbf P^1}C_U\). In particular \(Y\to X\)
is Q-Galois and \(Y\to C_U\) has degree n. These are the
actual maps of HD-SS11, not newly assigned sources.

Above an original branch value with Q-inertia \(q_d\), a point
of X having local degree c has normalized pullback inertia
conjugate to \(q_d^c\), of order \(d/\gcd(d,c)\).
For the three prime orders, every genuine descended inertia is
one of the powers covered by (SAT.13), up to conjugacy. Original
marks killed in Q, and points above unramified values of the
Q-cover, create no genuine Q-inertia. Thus the complete descended
rank sum \(D'\), with no omitted positive-cohomology term, obeys
\[
 \begin{gathered}
 D'=6+\mu-2b(h-1),\qquad
 \sum_j(1-1/o_j)
       \le\frac{180}{7b}\sum_j\operatorname{rank}(1-U(q'_j))
       =\frac{180}{7b}D'.
 \end{gathered}
 \tag{SAT.14}
\]
Here each genuine \(o_j\) is 2, 3 or 7, its weight is at most
\(6/7\), and its rank is at least \(b/30\). No bound on the
number of descended branch points is being assumed.

\paragraph{The numerical source bound and exact-field input.}
The actual Q-cover of the sphere has excess \(1/42\). Apply
Riemann--Hurwitz first to the degree-n map \(Y\to C_U\) and
then to \(Y\to X\), and substitute (SAT.14). This gives
\[
 \begin{split}
 \frac n{42}&\le 2h-2+\frac{180}{7b}
                         \{6+\mu-2b(h-1)\},\\
 n&\le2076-2076h+\frac{1080\mu}{b}+\frac{6480}{b}
       \le3156+\frac{4320}{b}\le4596.
 \end{split}
 \tag{SAT.15}
\]
The exact imaginary-quadratic K-model and transitivity imply
\(\mu\le b-\varepsilon_K(b)\le b-2\) by HD-SS09,
(SH.8). Thus \(\mu>0\) forces \(b\ge3\), and \(h\ge0\)
justifies all the displayed inequalities. This is precisely where
the actual field assumption is retained; no larger coefficient
field is substituted. This proves (SAT.5).

\paragraph{The original complete decomposition and minimality.}
In the stated minimal setting, D-TD44 gives transitivity and
\(\mu>0\) with the full original complement. It also gives
\[
 n\le108\max\{80,m_*-1\}-12,\qquad
 n>8628\ \Longrightarrow\
          m_*\ge\left\lceil\frac{n+120}{108}\right\rceil\ge82.
 \tag{SAT.16}
\]
Here \(m_*\) is the largest alternating natural degree among all
remaining components of that one fixed complete decomposition,
including permuted and nonstandard components. It is not a newly
chosen recursion of an unaccounted partner. If \(n>8628\), choose
an actual coarse layer containing that degree. It satisfies
(SAT.3), with \(m=m_*\ge82\) and \(mt\ge82>27\).
Its module need not be standard or stable componentwise. The
already proved (SAT.5) would give \(n\le4596\), a contradiction.
This proves (SAT.6).

The minimality input in D-TD44 is exactly its previously checked
post-vanishing actual-source descent: \(\mu=0\) would preserve
the exact K-model, signature, rank-nine variation and natural
signed polarization determinant while lowering coefficient
degree. We neither repeat that descent with positive \(\mu\)
nor assign the sixfold budget to U itself.
\end{proof}

\paragraph{Scope and dependency boundary.}
The new argument combines the actual complete extraction interface,
the support-sensitive conjugate bound D-TD30--31, the quarter-support
tensor-permutation estimate D-TD14, and the actual connected-cover
and exact-field identities of HD-SS08--11. The root contribution
and the separate D-TD47 generalization remain explicitly recorded.
There is no new external theorem, character classification, field
descent or period-dominance premise. The first assertion is a
triangle-source obstruction with one actual alternating layer;
the second uses D-TD44 to force such a layer. Other Q-signatures,
all nonminimal realizations outside the first assertion, and a
global source-degree cutoff remain open. This does not classify
the next large layer as a standard representation or iterate a
rank-nine source budget on a positive-cohomology partner.
